% Generated draft: requires mathematical and source review.
\documentclass[11pt]{article}
\usepackage[margin=25mm]{geometry}
\usepackage{fontspec}
\setmainfont{Latin Modern Roman}
\usepackage{amsmath,amssymb,mathtools,unicode-math}
\setmathfont{Latin Modern Math}
\usepackage{enumitem,needspace,xurl,hyperref}
\setlength{\parindent}{0pt}
\setlength{\parskip}{5pt}
\setlength{\emergencystretch}{3em}
\newcommand{\N}{\mathbb{N}}
\newcommand{\Z}{\mathbb{Z}}
\newcommand{\Q}{\mathbb{Q}}
\newcommand{\R}{\mathbb{R}}
\newcommand{\C}{\mathbb{C}}
\newcommand{\F}{\mathbb{F}}
\newcommand{\A}{\mathbb{A}}
\newcommand{\PP}{\mathbb{P}}
\newcommand{\E}{\mathbb{E}}
\newcommand{\eps}{\varepsilon}
\newcommand{\dd}{\,\mathrm{d}}
\begin{document}
\section*{UnsolvedMath ID: 10400160}
\textbf{Source problem code:} catalogue entry
\subsection{Definitions and mathematical statement}
Kerler's length-one TQFT programme has four stated tasks, using the source's length \(L\), mapping-class groups \(\Gamma_g\), algebraic object \(Q^*\) and theories \(V_p^{[0]}\):
\begin{enumerate}
\item Describe and construct algebraic \(L=1\) extensions of \(\Gamma_g\)-representations to TQFTs, preferably as simple generalizations of the Frohman--Nicas \(U(1)\) theory.
\item Classify \(L=1\) TQFTs through an extension theory of \(Q^*\).
\item Identify the \(\Gamma_g\)-representations on relative \(SU(2)\) moduli spaces from [77] with these TQFTs, and find analogous higher-rank theories.
\item If possible, identify \(V_p^{[0]}\) as length-one theories.
\end{enumerate}
The notions of length, extension, \(Q^*\), relative moduli spaces and \(V_p^{[0]}\) are not defined in this extract. Thus these are construction/classification targets with unresolved input conventions, not four precise existence theorems with invented hypotheses. Suggested generalizations in (1) use other, possibly noncompact gauge groups, or refined intersection homology of stratified moduli spaces. Given a theory for \(Q^*\), the source says (2) leads to problems in \(\mathfrak{sp}\)-invariant theory. It reports (4) for \(p=5\) in [211].

The following exposition refines finite length to \(q/l\)-solvability [212]: a TQFT over \(B[y]/(y^{l+1})\), whose constant-order theory over \(B\) has length \(q\). Here \(B\) renames the source's ground ring \(\mathrm M\) to distinguish it from a manifold. Such a theory is stated to have length at most \(ql+q+l\). The source reformulates Murakami [291] as producing from Reshetikhin--Turaev theory a \(1/1\)-solvable theory \(V_p^{[\leq1]}\), over \(\F_p[y]/(y^2)\), described as length three, with
\[V_p^{[\leq1]}(M)=1+\frac{y}{6}\lambda_{\mathrm{CWL}}(M)\tag{47}\]
for closed homology spheres \(M\). The meaning of division by six and reduction of \(\lambda_{\mathrm{CWL}}\) requires an admissible prime and coefficient convention. It reports an extension to rational homology spheres following Ohtsuki, with additional renormalization involving \(|H_1(M)|\), and calls a theory with this property of Casson type. It also says that for \(b_1(M)\geq1\), the related formula (46) for \(\lambda_{\mathrm{CWL}}\) is already encoded by a homological \(L=0\) TQFT through a special evaluation of Turaev--Milnor torsion [212]. Formula (46) is not reproduced.

For an admissible pair \((M,\varphi)\), with \(\varphi:\pi_1(M)\twoheadrightarrow\Z\), let \(\Sigma\) be a surface dual to \(\varphi\), and let \(C_\Sigma=M\setminus\Sigma:\Sigma\to\Sigma\) denote the cut cobordism. The source describes
\[V(M,\varphi)=\operatorname{tr}(A),\qquad A=V(C_\Sigma),\]
and a Turaev--Viro module, up to conjugacy,
\[M_{\mathrm{TV}}(M,\varphi)=\bigl(V(\Sigma)/\ker A^N,\ \overline A\bigr)\quad(N\text{ sufficiently large}),\]
with the induced endomorphism \(\overline A\). It gives traces of powers and dimension as numerical invariants. Eventual stabilization of the kernels and the meaning of dimension require finiteness/ground-ring hypotheses not supplied here. The Frohman--Nicas theories \(V^{(j)}\) are said to give Alexander-polynomial coefficients and hence \(\lambda_{\mathrm{CWL}}\) [212]; the module is attributed to Gilmer's description [143] as a refinement roughly generalizing the Alexander module.

For a \(1/1\)-solvable theory the source writes uniquely
\[V(M,\varphi)=\lambda^V_\varphi(M)+y\,\Xi^V_\varphi(M)\in B[y]/(y^2).\]
It states that if \(y\) is the half-projective parameter, \(\lambda^V_\varphi\) is independent of \(\varphi\), and expects it to be a function of \(\lambda_{\mathrm{CWL}}\). It further says that descent from a \(1/2\)-solvable TQFT with the same property would also make \(\Xi^V_\varphi\) independent of \(\varphi\). The half-projective parameter and descent are undefined here. For the Reshetikhin--Turaev-derived theory over \(\F_5[y]/(y^2)\), it reports this invariant defined in [212], and suggests an integral lift \(\Xi_{\Z}\). For \(p>5\), it expects next-order terms in the expansions (46) to be polynomial expressions in \(\lambda_{\mathrm{CWL}}\) and \(\Xi_{\Z}\). These proposed refinements, lift and polynomial relations are expectations, not established consequences or specified formulas in the supplied excerpt.
\subsection{Short English statement}
Kerler asks for constructions and an extension-theoretic classification of length-one TQFTs, including moduli-space and modular examples. The surrounding text motivates further solvable theories and possible refinements of Casson-type and torsion invariants; its conditional claims and expectations are kept distinct.
\subsection{Sources}
\begin{itemize}
\item Local UnsolvedMath registry, permanent ID \(10400160\), exact source code \texttt{catalogue entry}.
\item T. Kerler, attribution in Ohtsuki, \emph{Problems on invariants of knots and 3-manifolds} (2002), Problem 8.14, PDF page 128, equation (47) and reference to (46): \url{https://msp.org/gtm/2002/04/gtm-2002-04-024s.pdf}. Registry access 2026-07-29, pdf-font-structure extraction.
\item Internal references [77] (relative \(SU(2)\) moduli), [211] (\(p=5\) identification), [212] (solvability/refinements), Murakami [291] and Gilmer [143]. Their full bibliographic entries are not supplied. Frohman--Nicas, Ohtsuki and Reshetikhin--Turaev constructions are named in the source exposition.
\end{itemize}
\end{document}
